\documentclass[11pt,a4paper]{ctexart}
\usepackage[margin=2.2cm]{geometry}
\usepackage{hyperref}
\usepackage{enumitem}
\usepackage{xcolor}
\usepackage{longtable}
\usepackage{array}
\setlist[itemize]{leftmargin=1.8em}
\setlist[enumerate]{leftmargin=2em}
\hypersetup{
  colorlinks=true,
  linkcolor=blue,
  urlcolor=blue
}

\title{GeoAnnotate 项目结构详细说明}
\author{}
\date{\today}

\begin{document}
\maketitle

\section{项目当前整体架构}

当前项目不是传统意义上的纯前后端分仓，而是一个单仓库混合结构：

\begin{itemize}
  \item 浏览器端页面与交互使用 Next.js
  \item 服务端逻辑主要使用 CloudBase 云函数
  \item 数据结构和数据库模式由 Prisma 管理
\end{itemize}

从职责上看，可以拆成四层：

\begin{enumerate}
  \item 前端层
  \item 后端层
  \item 前后端连接层
  \item 共享配置层
\end{enumerate}

\section{前端层}

前端层主要负责：

\begin{itemize}
  \item 页面展示
  \item 用户交互
  \item 标注框绘制
  \item 地图展示
  \item 登录与注册
  \item 管理后台页面
  \item Wiki 页面壳与导航
\end{itemize}

\subsection{主要目录}

\begin{itemize}
  \item \texttt{app/}
  \item \texttt{components/}
  \item \texttt{app/globals.css}
  \item \texttt{tailwind.config.ts}
  \item \texttt{next.config.js}
\end{itemize}

\subsection{页面路由结构}

\textbf{根页面}

\begin{itemize}
  \item \texttt{app/page.tsx}：项目落地页，用于区分 \texttt{/wiki} 和 \texttt{/app}
\end{itemize}

\textbf{核心应用页面}

\begin{itemize}
  \item \texttt{app/(main)/}：核心平台页面的主要实现位置
\end{itemize}

\textbf{认证页面}

\begin{itemize}
  \item \texttt{app/(auth)/}：登录、注册、找回密码
\end{itemize}

\textbf{管理后台}

\begin{itemize}
  \item \texttt{app/admin/}：管理端页面
\end{itemize}

\textbf{新路由分区}

\begin{itemize}
  \item \texttt{app/app/}
  \item \texttt{app/auth/}
\end{itemize}

这两个目录更多是新的 URL 入口层，很多文件只是转发到原有页面实现。

因此后续真正要改页面内容时，优先修改：

\begin{itemize}
  \item \texttt{app/(main)/...}
  \item \texttt{app/(auth)/...}
\end{itemize}

\subsection{前端组件目录}

\begin{itemize}
  \item \texttt{components/layout/}：导航栏、整体页面壳、路由壳
  \item \texttt{components/annotation/}：标注相关组件
  \item \texttt{components/map/}：地图显示与选择
  \item \texttt{components/battle/}：计时器、计分板
  \item \texttt{components/wiki/}：Wiki 导航
  \item \texttt{components/ui/}：通用 UI 组件
\end{itemize}

\subsection{前端常见修改位置}

\textbf{改首页}

\begin{itemize}
  \item \texttt{app/page.tsx}
  \item \texttt{app/(main)/home/page.tsx}
\end{itemize}

\textbf{改标注相关界面}

\begin{itemize}
  \item \texttt{app/(main)/annotate/mode/page.tsx}
  \item \texttt{app/(main)/annotate/page.tsx}
  \item \texttt{components/annotation/BBoxCanvas.tsx}
  \item \texttt{components/annotation/ThoughtInput.tsx}
  \item \texttt{components/map/LocationMap.tsx}
\end{itemize}

\textbf{改 AI 对战界面}

\begin{itemize}
  \item \texttt{app/(main)/battle/page.tsx}
  \item \texttt{app/(main)/battle/[sessionId]/play/page.tsx}
  \item \texttt{app/(main)/battle/[sessionId]/result/page.tsx}
  \item \texttt{components/battle/CountdownTimer.tsx}
  \item \texttt{components/battle/BattleScoreBoard.tsx}
\end{itemize}

\textbf{改整体布局}

\begin{itemize}
  \item \texttt{components/layout/Navbar.tsx}
  \item \texttt{components/layout/AppShell.tsx}
  \item \texttt{components/layout/MainRouteShell.tsx}
  \item \texttt{app/admin/layout.tsx}
  \item \texttt{app/(main)/wiki/layout.tsx}
\end{itemize}

\section{后端层}

当前项目的后端不是独立 NestJS / FastAPI 服务，而是以 CloudBase 云函数作为主要后端执行层。

\subsection{主要目录}

\begin{itemize}
  \item \texttt{cloudfunctions/}
  \item \texttt{prisma/}
\end{itemize}

\subsection{后端负责的内容}

\begin{itemize}
  \item 任务分发
  \item 标注提交
  \item 用户信息读取
  \item AI 对战创建
  \item AI 对战结果读取
  \item 奖励与积分逻辑
  \item 排行榜与历史数据读取
  \item 管理后台数据处理
\end{itemize}

\subsection{关键云函数}

\begin{itemize}
  \item \texttt{cloudfunctions/get-next-task/}
  \item \texttt{cloudfunctions/submit-annotation/}
  \item \texttt{cloudfunctions/create-battle/}
  \item \texttt{cloudfunctions/submit-battle-round/}
  \item \texttt{cloudfunctions/get-battle-result/}
  \item \texttt{cloudfunctions/get-user-profile/}
  \item \texttt{cloudfunctions/get-user-history/}
  \item \texttt{cloudfunctions/get-points-history/}
  \item \texttt{cloudfunctions/get-leaderboard/}
  \item \texttt{cloudfunctions/list-images/}
  \item \texttt{cloudfunctions/review-annotation/}
  \item \texttt{cloudfunctions/export-annotations/}
\end{itemize}

\subsection{数据库层}

数据库结构主要在：

\begin{itemize}
  \item \texttt{prisma/schema.prisma}
\end{itemize}

如果要改数据库结构，优先修改这个文件，并同步检查对应云函数和前端调用层。

\section{前后端连接层}

前端和后端之间并不是页面直接访问数据库，而是通过统一的调用层转发到云函数。

最关键的文件是：

\begin{itemize}
  \item \texttt{lib/cloudbase.ts}
\end{itemize}

它的作用包括：

\begin{itemize}
  \item 封装 CloudBase SDK 初始化
  \item 封装登录、注册、找回密码
  \item 封装云函数调用
  \item 给前端页面提供统一的方法名
\end{itemize}

如果后端新增或修改了接口，通常需要同步修改：

\begin{itemize}
  \item \texttt{cloudfunctions/...}
  \item \texttt{lib/cloudbase.ts}
\end{itemize}

\section{配置与共享层}

这部分不直接属于前端界面，也不属于后端业务逻辑，但它决定了项目的整体骨架。

\subsection{主要目录和文件}

\begin{itemize}
  \item \texttt{features/}
  \item \texttt{types/}
  \item \texttt{lib/modes.ts}
  \item \texttt{package.json}
  \item \texttt{tsconfig.json}
  \item \texttt{next.config.js}
\end{itemize}

\subsection{各部分作用}

\begin{itemize}
  \item \texttt{features/annotation/}：标注模式与标注类型配置
  \item \texttt{features/battle/}：对战模式、AI 对手、回合数、时间配置
  \item \texttt{features/admin/}：管理后台导航配置
  \item \texttt{features/platform/}：平台入口与路由结构说明
  \item \texttt{types/}：共享类型定义
  \item \texttt{lib/modes.ts}：统一导出模式配置
\end{itemize}

\section{如何快速判断改哪里}

\begin{longtable}{|>{\raggedright\arraybackslash}p{0.22\textwidth}|>{\raggedright\arraybackslash}p{0.32\textwidth}|>{\raggedright\arraybackslash}p{0.34\textwidth}|}
\hline
\textbf{目标} & \textbf{优先修改位置} & \textbf{说明} \\
\hline
改页面样式和布局 & \texttt{app/(main)}、\texttt{components} & 前端页面与组件 \\
\hline
改导航和壳层 & \texttt{components/layout}、\texttt{app/admin/layout.tsx} & 页面骨架与导航结构 \\
\hline
改接口调用方式 & \texttt{lib/cloudbase.ts} & 前后端调用桥梁 \\
\hline
改业务逻辑 & \texttt{cloudfunctions} & 云函数处理数据库和逻辑 \\
\hline
改数据库结构 & \texttt{prisma/schema.prisma} & 数据表与字段定义 \\
\hline
改模式与配置 & \texttt{features}、\texttt{lib/modes.ts} & 全局配置与共享常量 \\
\hline
\end{longtable}

\section{一句话总结}

当前项目可以简单理解为：

\begin{itemize}
  \item \texttt{app/ + components/} = 前端
  \item \texttt{cloudfunctions/ + prisma/} = 后端
  \item \texttt{lib/cloudbase.ts} = 前后端连接层
  \item \texttt{features/ + types/ + 配置文件} = 项目框架与共享配置
\end{itemize}

如果后续要同步到 GitHub README，建议用精简版放在 README 中，把详细版长期维护在 \texttt{docs/} 目录。

\end{document}
